\section{Emergent Abilities：能力跳变还是测量产物}

课堂首先回顾 2022 年关于大模型涌现能力的经典定义：某项能力在小模型上接近随机，达到某个规模后才明显出现，而且不能由小模型趋势直接外推。但讲者也立即补上测量争议：离散指标、prompt 和任务定义都可能改变曲线形状。因此，正确问题不是“涌现是真是假”的二选一，而是哪些行为、在哪种测量下、为何表现为阈值。

\subsection{经典定义与 critical threshold}

第一张标题页把 emergence 作为独立研究对象，随后定义页给出三个条件：大模型有而小模型没有、不能由小模型直接外推、性能在 critical threshold 前近随机而后快速上升。这个定义强调 observed behavior，不提供内部机制解释。本节先固定这套历史定义，后面再用 prompt 与 metric 争议检验它的证据边界。

\lecturefigure{slide-02-emergent-title.jpg}{Emergent Abilities of Large Language Models}{Stanford 课堂视频 00:00:33}

\begin{knowledgebox}{读图：标题页的作用是建立研究对象}
这里的 ability 可以是 modular arithmetic、word unscrambling、question answering 或 prompting behavior。它不是单个 neuron 突然“学会技能”，而是 evaluation pipeline 观察到的任务性能模式；模型、prompt、scoring 和 task distribution 都属于现象定义的一部分。
\end{knowledgebox}

\lecturefigure{slide-03-emergent-definition.jpg}{课堂采用的涌现能力定义}{Stanford 课堂视频 00:00:39}

\begin{importantbox}{阈值定义的数学直觉}
设规模为 $s$，任务得分为 $m(s)$。经典涌现叙事认为存在 $s_c$，使
\[
m(s)\approx m_{\text{random}}\quad(s<s_c),\qquad
m(s)\gg m_{\text{random}}\quad(s\ge s_c).
\]
$s_c$ 是观察到的 critical scale，$m_{\text{random}}$ 是随机基线。这个式子描述曲线形状，不证明参数空间发生真正物理相变。
\end{importantbox}

\teachervoice{讲者用 phase transition 类比规模阈值，但同时说研究者至今仍缺少充分解释。课堂的“不可预测”是指从已测小模型曲线做朴素外推会失败，不是说任何能力都完全无法研究或预测。}

\subsection{Few-shot 曲线怎样读}

多任务图把 training FLOPs 放在横轴，把任务 performance 放在纵轴；LaMDA、GPT-3、Gopher、Chinchilla、PaLM 等模型在 modular arithmetic、IPA transliteration、word unscrambling 和 Persian QA 上出现不同阈值。读图时不能只看“最大模型最高”，还要看小模型区间是否平坦、任务是否采用 exact match、不同模型族是否落在同一趋势上。

\lecturefigure{slide-04-few-shot-emergence.jpg}{多项 few-shot 任务随 training FLOPs 变化的曲线}{Stanford 课堂视频 00:01:18}

\begin{knowledgebox}{读图：先确认横纵轴，再判断跳变}
横轴是训练计算量的对数尺度，纵轴因任务而异。某些图在 $10^{22}$--$10^{23}$ FLOPs 左右明显上升，但点数有限且模型架构不同。若 metric 是全对或全错，连续的 token-level 改善可能被压成突然 crossing，因此阈值既可能包含真实能力变化，也可能包含 scoring nonlinearity。
\end{knowledgebox}

\begin{warningbox}{不要从少数离散点推导必然相变}
模型规模、数据、架构和训练 recipe 同时变化；曲线上的点不是只改变一个自变量的 controlled experiment。涌现图支持“朴素小模型外推不可靠”，不自动支持“内部表征在某一点突然出现”。
\end{warningbox}

\subsection{Prompting strategy 会改变所观察到的能力}

下一组图展示 augmented prompting，包括 scratchpad、instruction following、八位数加法和 calibration。最典型例子是 GSM8K 上，普通 prompting 随规模增长有限，而 Chain-of-Thought 在大模型上显著提高。这里的能力既来自参数，也来自把计算过程暴露在输出 token 中。

\lecturefigure{slide-05-augmented-prompting.jpg}{Augmented prompting 下的规模与能力变化}{Stanford 课堂视频 00:01:54}

\begin{knowledgebox}{读图：模型能力与 elicitation method 不能分离}
同一 checkpoint 使用不同 prompt 可得到完全不同曲线。若 CoT 让大模型把隐含多步计算转换成显式序列，观测到的“涌现”可能是 latent capability 与 elicitation threshold 的组合。评估报告因此必须同时写 model、prompt examples、decoding 和 metric。
\end{knowledgebox}

\begin{importantbox}{观测性能是组合函数}
更完整地写，任务得分是
\[
m=G(\theta, p, d, \mu),
\]
其中 $\theta$ 表示模型参数与训练，$p$ 表示 prompt/inference strategy，$d$ 表示 decoding，$\mu$ 表示 metric。只把 $m$ 对参数规模作图，会把其他三项的非线性一起投影到“规模效应”中。
\end{importantbox}

\subsection{能力目录与规模只是索引}

课堂保留论文中的大表，列出不同 emergent abilities、对应 training FLOPs、参数量、模型和参考文献。它的价值是让研究者发现跨任务模式，而不是给每种能力一个永久固定的“最小参数量”。因此本节把表格当作实验索引，重点检查任务、模型族和评测设置，而不是背诵门槛数字。

\lecturefigure{slide-06-emergent-abilities-table.jpg}{论文汇总的涌现能力与出现规模}{Stanford 课堂视频 00:02:33}

\begin{knowledgebox}{读图：表格中的 scale 不是普适门槛}
相同任务在不同 tokenizer、data mixture、instruction tuning 和 prompt 下会移动阈值；不同模型族的参数利用率也不同。先看 task 和 reference，再看 FLOPs/parameters。规模列更像历史实验索引，而不是设计新模型时可直接查表的常数。
\end{knowledgebox}

\subsection{为什么会涌现仍未解决}

解释页明确说，目前对原因了解有限；使用的 evaluation metrics 也可能无法充分解释现象，并建议用 cross-entropy loss 等连续指标分析。连续 metric 能揭示小模型已经发生但未跨过 exact-match 门槛的渐进改善。本节从“看到跳点”转向“解释跳点”，把 measurement artifact 与能力机制分开讨论。

\lecturefigure{slide-07-emergence-explanations.jpg}{涌现原因与 evaluation metric 的不确定性}{Stanford 课堂视频 00:02:45}

\begin{knowledgebox}{读图：metric critique 改变的是证据强度}
若 exact match 只有完整答案正确才记 1 分，模型从“几乎正确”到“正确”会产生跳点；cross-entropy 则能显示概率质量逐步迁移。连续曲线可能削弱“突然相变”的视觉效果，但仍不能否定复杂行为在某个实用尺度后才可用。
\end{knowledgebox}

\begin{warningbox}{Mirage critique 也不是“能力没有变化”}
指标导致 discontinuity，并不等于规模无效；它说明“能力突然出现”与“得分突然跨线”必须分开。工程上仍需关心某项能力何时达到可用阈值，只是不能把阈值当作唯一机制证据。
\end{warningbox}

\teachervoice{讲者没有给出单一 emergence theory，而是把 metric 选择列为核心研究问题。这个谨慎边界比背下阈值更重要：先问现象怎样测，再问模型内部发生了什么。}

\subsection{Beyond scaling 的研究路线}

即使扩大规模可能带来新能力，课堂也强调它不是唯一变量。新架构、更高质量数据、更好的 training procedure、更强 few-shot prompting、理论与 interpretability、computational linguistics 都可能让能力更早出现或更可控。

\lecturefigure{slide-08-beyond-scaling.jpg}{除继续放大之外的能力研究方向}{Stanford 课堂视频 00:03:06}

\begin{knowledgebox}{读图：这些方向作用于不同层}
Architecture 和 training 改变参数学习；data quality 改变监督信号；prompting 改变 inference interface；interpretability 与 theory 改变解释和预测能力；computational linguistics 提供结构先验。把它们都压成“有效算力”会失去可干预性。
\end{knowledgebox}

\subsection{能力增长伴随风险增长}

Emergent risk slide 把 truthfulness、bias、toxicity、memorization 和未知风险放在一起。更大模型能更好完成任务，也可能更有效重现敏感数据、适应有害提示或在含糊语境中放大偏差；能力与风险不共享一个单调安全轴。因此本节把安全评测视为独立维度，而不是从平均任务分数推断风险已经下降。

\lecturefigure{slide-09-emergent-risks.jpg}{规模增长可能同时放大能力与社会风险}{Stanford 课堂视频 00:04:00}

\begin{knowledgebox}{读图：风险不是训练后再加的附录}
Truthfulness 是事实与校准问题，bias 涉及群体差异，toxicity 涉及有害输出，memorization 涉及隐私与版权。它们需要不同 dataset、attack 和 metric。一个总体 helpfulness 分数不能证明这些风险同步下降。
\end{knowledgebox}

\teachervoice{讲者提醒，未来模型还可能出现尚未发现的风险。这里的 unknown risk 不是制造恐慌，而是说明 capability discovery 与 safety discovery 都需要持续 red teaming 和 measurement。}

\subsection{General-purpose model 改变研究与应用结构}

Sociological changes slide 讨论的不只是技术扩散。通用模型让社区从“每个任务训练一个模型”转向 shared foundation，再通过 in-context prompting 或轻量适配服务多个领域；Transformer 也从 NLP 扩展到 text-to-image、biology、robotics 等场景。

\lecturefigure{slide-10-sociological-changes.jpg}{涌现与通用模型带来的社区和应用变化}{Stanford 课堂视频 00:04:45}

\begin{knowledgebox}{读图：技术范式会重分配进入门槛}
Shared model 降低应用开发成本，却把训练权力、评估接口和基础设施集中到少数提供者。研究者可以更快调用能力，也更依赖不可见 data/updates。BabyLM 稍后正是对这种资源集中提出反问题。
\end{knowledgebox}

\subsection{未来工作与课堂开放问题}

Future-work slide 把规模、架构、训练、数据、prompting、小模型能力、theory 与 computational linguistics 汇总；讨论页进一步追问规模是否有极限、数据与 metric 在 emergence 中扮演什么角色，以及 specialized systems 是否可能优于“把一切塞进参数”的 monolith。

\lecturefigure{slide-11-future-work.jpg}{Emergent abilities 的未来研究方向}{Stanford 课堂视频 00:06:00}

\begin{knowledgebox}{读图：future work 不是单一 scaling roadmap}
列表同时包含 capability、efficiency、interpretability 和 linguistic structure。最值得保留的是多路线比较：同一能力可以通过更多参数、更好数据、更好 prompt、外部工具或专门模型获得，研究目标应比较成本与可靠性，而不只比较 peak score。
\end{knowledgebox}

\lecturefigure{slide-12-discussion-questions.jpg}{讲者留给课堂的 emergence 开放问题}{Stanford 课堂视频 00:06:57}

\begin{knowledgebox}{读图：问题本身就是证据边界}
这组问题没有标准答案：能力会否持续出现、规模是否递减、metric 与 data 是否制造阈值、专业系统是否更可靠。讲义不应把 2026 观点倒灌成讲者结论，而应把这些问题作为后续 reasoning、BabyLM 和 agents 三个单元的连接点。
\end{knowledgebox}

\teachervoice{讲者邀请学生随时打断并讨论，说明这部分不是宣布 emergence 定律，而是用经典论文建立共同问题空间。}

\subsection{本章小结}

经典涌现定义描述阈值式任务得分，但曲线由模型、prompt、decoding 和 metric 共同决定。规模可能带来新能力，也带来新风险和资源集中；更好的数据、架构、推理程序与专门系统都是替代或互补路线。

